\section{问题建模与自适应卡尔曼滤波}

\subsection{状态空间模型}

考虑车辆在局部平面内的运动，选取状态向量为 $\mathbf{x}_k = [x_k, y_k, v_{x,k}, v_{y,k}, a_{x,k}, a_{y,k}]^\top$，其中 $x_k$ 与 $y_k$ 表示横纵向位置，$v_{x,k}$ 与 $v_{y,k}$ 为对应速度分量，$a_{x,k}$ 与 $a_{y,k}$ 为加速度。离散时间状态转移方程采用匀加速模型：
\begin{equation}
    \mathbf{x}_{k+1} = \mathbf{F} \mathbf{x}_k + \mathbf{\Gamma} \mathbf{w}_k,
\end{equation}
其中 $\mathbf{F} \in \mathbb{R}^{6\times 6}$ 为状态转移矩阵，$\mathbf{\Gamma} \in \mathbb{R}^{6\times 3}$ 为噪声驱动矩阵，过程噪声 $\mathbf{w}_k \sim \mathcal{N}(\mathbf{0}, \mathbf{Q}_k)$。观测模型为
\begin{equation}
    \mathbf{z}_k = \mathbf{H} \mathbf{x}_k + \mathbf{v}_k,
\end{equation}
$\mathbf{z}_k$ 通常包含位置与速度观测量，$\mathbf{v}_k \sim \mathcal{N}(\mathbf{0}, \mathbf{R}_k)$ 为观测噪声。

\subsection{自适应噪声协方差估计}

传统卡尔曼滤波假设 $\mathbf{Q}_k$ 与 $\mathbf{R}_k$ 是先验已知的常量，这在实际场景中往往难以满足。本文基于滑动窗口内新息序列的统计信息，对过程噪声协方差 $\mathbf{Q}_k$ 进行在线估计：
\begin{equation}
    \hat{\mathbf{Q}}_k = \frac{1}{W} \sum_{i=k-W+1}^{k} \mathbf{K}_{i-1} \mathbf{r}_i \mathbf{r}_i^\top \mathbf{K}_{i-1}^\top,
\end{equation}
其中 $W$ 为窗口长度（选取 $W=30$），$\mathbf{r}_i$ 为第 $i$ 时刻的新息向量。观测噪声协方差 $\mathbf{R}_k$ 通过稳健的M估计器根据残差调整，以抵御异常值影响。

\subsection{异常检测与野值剔除}

在城市遮挡环境下，GPS或视觉定位常出现瞬时跳变。本文采用基于马氏距离的决策规则：若新息满足 $\mathbf{r}_k^\top \mathbf{S}_k^{-1} \mathbf{r}_k > \tau$，则将当前观测视为野值并跳过更新步骤；阈值 $\tau$ 根据自由度为6的卡方分布设定在置信水平95%处。该机制在保证正常观测精度不降低的前提下，有效提高了滤波器对局部遮挡的鲁棒性。